\documentclass[10pt,a4paper]{amsart}
\usepackage[margin=19mm]{geometry}
\usepackage{amsmath,amssymb,mathtools}
\usepackage[scr=rsfs,cal=euler]{mathalfa}
\usepackage{xfrac}
\usepackage[unicode,hidelinks,bookmarks=false]{hyperref}
\newcommand{\ie}{i.e.\ }
\newcommand{\cf}{cf.\ }
\newcommand{\resp}{respectively\ }
\newcommand{\Subsection}{\subsection}
\providecommand{\nobreakdash}{\nobreak\hbox{-}}
\setlength{\emergencystretch}{2em}
\multlinegap0pt
\usepackage{bm}
\usepackage{bbm}
\usepackage{dsfont}
\usepackage{xspace}
\usepackage{cancel}
\usepackage{mathtools}
\usepackage{amsthm}
\makeatletter
\@ifundefined{prop}{\newtheorem{prop}{Prop}}{}
\makeatother
\DeclareMathOperator{\diam}{diam}
\newcommand{\I}{\int\limits}
\newcommand{\R}{\mathbb{R}}
\title[$C^{1,\alpha}$-regularity for Schr\"odinger potentials of Sh]{$C^{1,\alpha}$-regularity for Schr\"odinger potentials of Shannon-entropy regularized Optimal Transport}
\author{Sumiya Baasandorj, Simone Di Marino, Augusto Gerolin}
\date{Source: arXiv:2506.09302v1 (CC BY 4.0)}
\begin{document}
\maketitle

\noindent\emph{Source and licence.} $C^{1,\alpha}$-regularity for Schr\"odinger potentials of Shannon-entropy regularized Optimal Transport. Version arXiv:2506.09302v1 (CC BY 4.0), distributed under the Creative Commons Attribution 4.0 International licence, as recorded in the arXiv licence metadata for this exact version and shown on the abstract page https://arxiv.org/abs/2506.09302v1. Full rights notice: licenses/arxiv-2506.09302-CC-BY-4.0.txt.\par\smallskip

\noindent\textit{Editorial scope.} Counted unit: the Prop whose source label is \texttt{prop:convex} in the article (source0.tex lines 359--369, source label \texttt{prop:convex}), delivered together with the article's own complete original proof of it (source0.tex lines 370--412, one \texttt{proof} environment of 43 lines). No mathematics is written, repaired or re-derived: statement and proof are the article's own text line for line. Required context retained uncounted: none. Every definition, display and label of those lines is retained unchanged. Omitted from this package and not redistributed: 1--358, 413--937. Nothing inside a retained excerpt is omitted or altered: every retained body is delivered exactly as the article's lines are written, so no omission receipt of any kind accompanies this package. Citations: the retained excerpts carry no citation command at all, so no bibliography entry of the article is transcribed and no reference is redistributed. Reference adaptation: none was needed and none was made; every reference inside the delivered unit resolves inside it, so no unresolved reference or citation is produced. Printed slips kept literal: no printed word, symbol, hypothesis or number of the counted statement or of its proof was altered. Layout and class: the article's own class is \texttt{article} with packages including amsfonts, amsmath, amssymb, amsthm, hyperref, color, dsfont, fontenc, authblk, graphicx, tikz, xcolor, eucal, enumerate; the wrapper is the lane's pinned \texttt{amsart} editorial wrapper at 10pt on A4, which restates the theorem-like environments the retained text needs and adds the article's own macro definitions that the retained text needs (\texttt{\textbackslash I}, \texttt{\textbackslash R}, \texttt{\textbackslash diam}), with extra wrapper packages (bm, bbm, dsfont, xspace, cancel, mathtools). The delivered numbering is the unit's own. Assets: no retained span contains a figure environment, a graphics-inclusion command, a PDF-page inclusion, a table or a TikZ picture; no image file is copied into this package, the package declares no asset dependency and no image file is redistributed with it. Licence: version arXiv:2506.09302v1 of this article is distributed under the Creative Commons Attribution 4.0 International licence, as recorded in the arXiv licence metadata for this exact version and shown on the abstract page https://arxiv.org/abs/2506.09302v1. A full rights notice is delivered at licenses/arxiv-2506.09302-CC-BY-4.0.txt. Characters: every retained excerpt is pure ASCII, so the delivered engine, XeLaTeX with its default Latin Modern text font, reports no missing character.
\begin{prop}
    \label{prop:convex}
    Let $\Omega\subset \R^{n}$ be a bounded open convex set. There exists a constant $c\equiv c(n,\Omega)$ such that 
    \begin{align*}
        %\label{prop:convex:1}
        \I_{\Omega\cap B_{r}(z)}\exp\left( -\frac{|x-z|}{r} \right)\,dx 
        \geqslant
        c|B_{r}(z)|
    \end{align*}
    holds whenever $z\in\overline{\Omega}$ and $0<r<1$.
\end{prop}

\begin{proof}
    First, we show that there exists a constant $c\equiv c(n,\Omega)$ such that 
    \begin{align}
        \label{convex:1}
        \frac{|\Omega\cap B_{r}(z)|}{|B_{r}(z)|} \geqslant c
    \end{align}
    for every ball $B_{r}(z)$ centered at $z\in \overline{\Omega}$ and of radius $0<r<1$. Let $B_{2R}(z_0)\subset \Omega$ be one of the largest balls contained in $\Omega$.  Let us consider a ball $B_{r}(z)$ centered at $z\in \overline{\Omega}$ and of radius $0<r<1$. Then we denote $Z_{R}$ a convex hull of $B_{R}(z_0)$ and $z$. We consider two cases depending on the position of center $z$.  

    \textbf{Case 1. $z\in \overline{\Omega}\setminus \overline{B_{R}}(z_0)$.} Let $B_{r_{z}}(x_{z})$ be the largest ball contained in $Z_{R}\cap B_{r}(z)$. $B_{r_z}(x_z)$ will be the image of $B_R(z_0)$ through an homotety centered at $z$ with ratio $\frac{r}{R+|z-z_0|}$;in particular we will have $r_z=\frac{ Rr}{R+|z-z_0|}$. In turn we get

     \begin{align*}
        \frac{|\Omega\cap B_{r}(z)|}{|B_{r}(z)|}
        \geqslant
        \frac{|B_{r_{z}}(x_{z})|}{|B_{r}(z)|}
        =
        \left( \frac{R}{R + |z-z_0|} \right)^{n}
        \geqslant 
         \left( \frac{R}{R + \diam(\Omega)} \right)^{n}
    \end{align*}

   \textbf{Case 2. $z\in \overline{B_{R}}(z_0)$. } For $0<r < R$, we have that $B_{r}(z)\subset B_{2R}(z_0)\subset \Omega$ and
 \begin{align*}
     \frac{|\Omega\cap B_{r}(z)|}{|B_{r}(z)|}
     \geqslant
     \frac{|B_{2R}(z_0)\cap B_{r}(z)|}{|B_{r}(z)|} \geqslant 1;
 \end{align*}
 for $1> r\geqslant R$, we find 
  \begin{align*}
     \frac{|\Omega\cap B_{r}(z)|}{|B_{r}(z)|}
     \geqslant
     \frac{|B_{2R}(z_0)\cap B_{r}(z)|}{|B_{r}(z)|} 
     \geqslant
     \frac{|B_{R}(0)|}{|B_{1}(0)|}
     \geqslant R^{n}.
 \end{align*}
Taking into account two cases we have discussed above, we arrive at the validity of \eqref{convex:1}. Finally, using inequality \eqref{convex:1}, we have 
\begin{align*}
    \I_{\Omega\cap B_{r}(z)} \exp\left(-\frac{|x-z|}{r}\right)\,dx
    \geqslant
    e^{-1}|\Omega\cap B_{r}(z)| \geqslant c |B_{r}(z)|
\end{align*}
for some constant $c\equiv c(n,\Omega)$. 
\end{proof}

\end{document}
